% Options for packages loaded elsewhere
\PassOptionsToPackage{unicode}{hyperref}
\PassOptionsToPackage{hyphens}{url}
\documentclass[
]{article}
\usepackage{xcolor}
\usepackage{amsmath,amssymb}
\setcounter{secnumdepth}{-\maxdimen} % remove section numbering
\usepackage{iftex}
\ifPDFTeX
  \usepackage[T1]{fontenc}
  \usepackage[utf8]{inputenc}
  \usepackage{textcomp} % provide euro and other symbols
\else % if luatex or xetex
  \usepackage{unicode-math} % this also loads fontspec
  \defaultfontfeatures{Scale=MatchLowercase}
  \defaultfontfeatures[\rmfamily]{Ligatures=TeX,Scale=1}
\fi
\usepackage{lmodern}
\ifPDFTeX\else
  % xetex/luatex font selection
\fi
% Use upquote if available, for straight quotes in verbatim environments
\IfFileExists{upquote.sty}{\usepackage{upquote}}{}
\IfFileExists{microtype.sty}{% use microtype if available
  \usepackage[]{microtype}
  \UseMicrotypeSet[protrusion]{basicmath} % disable protrusion for tt fonts
}{}
\makeatletter
\@ifundefined{KOMAClassName}{% if non-KOMA class
  \IfFileExists{parskip.sty}{%
    \usepackage{parskip}
  }{% else
    \setlength{\parindent}{0pt}
    \setlength{\parskip}{6pt plus 2pt minus 1pt}}
}{% if KOMA class
  \KOMAoptions{parskip=half}}
\makeatother
\usepackage{longtable,booktabs,array}
\newcounter{none} % for unnumbered tables
\usepackage{calc} % for calculating minipage widths
% Correct order of tables after \paragraph or \subparagraph
\usepackage{etoolbox}
\makeatletter
\patchcmd\longtable{\par}{\if@noskipsec\mbox{}\fi\par}{}{}
\makeatother
% Allow footnotes in longtable head/foot
\IfFileExists{footnotehyper.sty}{\usepackage{footnotehyper}}{\usepackage{footnote}}
\makesavenoteenv{longtable}
\usepackage{graphicx}
\makeatletter
\newsavebox\pandoc@box
\newcommand*\pandocbounded[1]{% scales image to fit in text height/width
  \sbox\pandoc@box{#1}%
  \Gscale@div\@tempa{\textheight}{\dimexpr\ht\pandoc@box+\dp\pandoc@box\relax}%
  \Gscale@div\@tempb{\linewidth}{\wd\pandoc@box}%
  \ifdim\@tempb\p@<\@tempa\p@\let\@tempa\@tempb\fi% select the smaller of both
  \ifdim\@tempa\p@<\p@\scalebox{\@tempa}{\usebox\pandoc@box}%
  \else\usebox{\pandoc@box}%
  \fi%
}
% Set default figure placement to htbp
\def\fps@figure{htbp}
\makeatother
\setlength{\emergencystretch}{3em} % prevent overfull lines
\providecommand{\tightlist}{%
  \setlength{\itemsep}{0pt}\setlength{\parskip}{0pt}}
\usepackage{bookmark}
\IfFileExists{xurl.sty}{\usepackage{xurl}}{} % add URL line breaks if available
\urlstyle{same}
\hypersetup{
  hidelinks,
  pdfcreator={LaTeX via pandoc}}

\author{}
\date{}

\begin{document}

\includegraphics[width=1.55in,height=0.82119in]{media/image1.png}

Assignment

STAT-4104

ID-12110070

Submitted to : Dr. Md. Siddikur Rahman ( Associate Professor)

\begin{figure}
\centering
\includegraphics[width=0.83in,height=0.52377in]{media/image2.png}
\caption{Submitted by : Khatune Jannat Rida}
\end{figure}

\textbf{\hfill\break
}

\textbf{Problem 1 :}

Considering the lung cancer data for solving problem 1 using python.

\textbf{(Basic EDA Question's answer )}

\textbf{Python code-}

Import Libraries and Dataset:

import pandas as pd

import numpy as np

from scipy.stats import chi2\_contingency, fisher\_exact

import statsmodels.formula.api as smf

from sklearn.metrics import roc\_auc\_score, confusion\_matrix

from sklearn.metrics import accuracy\_score, recall\_score

\# Load dataset

df = pd.read\_csv(\textquotesingle lung\_disease.csv\textquotesingle)

\# Display first 5 rows

print(df.head())

\textbf{i) Distribution of Age}

\textbf{Python Code :}

print(df{[}\textquotesingle Age\textquotesingle{]}.describe())

print(\textquotesingle Skewness =\textquotesingle,
df{[}\textquotesingle Age\textquotesingle{]}.skew())

\textbf{Output :}

count 500.000000

mean 43.904000

std 14.604841

min 20.000000

25\% 30.750000

50\% 45.000000

75\% 56.000000

max 69.000000

Skewness = 0.012

\textbf{Interpretation :}

 The average age is about 44 years.

 Minimum age = 20 and maximum age = 69.

 Skewness is very close to 0.

 Therefore, Age is approximately normally distributed.

 There is no strong positive or negative skew.

\textbf{ii) Proportion of Smokers vs Non-Smokers}

\textbf{Python Code :}

smoking\_prop =
df{[}\textquotesingle Smoking\textquotesingle{]}.value\_counts(normalize=True)
* 100

print(smoking\_prop)

\textbf{Output :}

No 58.6\%

Yes 41.4\%

\textbf{Interpretation :}

 About 41.4\% individuals are smokers.

 About 58.6\% are non-smokers.

 Non-smokers are more common in the dataset.

\textbf{iii) Income Group with Highest Frequency}

\textbf{Python Code :}

print(df{[}\textquotesingle Income\textquotesingle{]}.value\_counts())

\textbf{Output :}

Low 204

Medium 197

High 99

\textbf{Interpretation :}

 The Low income group has the highest frequency.

 The High income group has the lowest frequency.

\textbf{iv) Percentage Exposed to High Pollution}

\textbf{Python Code :}

high\_pollution = (df{[}\textquotesingle Pollution\textquotesingle{]} ==
\textquotesingle High\textquotesingle).mean() * 100

print(high\_pollution)

\textbf{Output :}

70.4\%

\textbf{Interpretation :}

 Around 70.4\% individuals are exposed to high pollution.

 High pollution exposure is very common in the dataset.

\textbf{v) Overall Prevalence of Lung Disease}

\textbf{Python Code :}

prevalence = (df{[}\textquotesingle LungDisease\textquotesingle{]} ==
\textquotesingle Yes\textquotesingle).mean() * 100

print(prevalence)

\textbf{Output :}

52.8\%

\textbf{Interpretation :}

 About 52.8\% individuals have lung disease.

 Lung disease prevalence is relatively high.

\textbf{(Association \& Crosstab)}

\textbf{i) Association Between Smoking and Lung Disease}

\textbf{Python Code :}

\begin{verbatim}
crosstab_smoking = pd.crosstab(df['Smoking'], df['LungDisease'])
print(crosstab_smoking)
\end{verbatim}

\textbf{Output :}

\begin{verbatim}
LungDisease   No   Yes
Smoking
No            184  109
Yes            52  155
\end{verbatim}

\textbf{Interpretation :}

 Among smokers, 155 individuals have lung disease.

 Among non-smokers, only 109 individuals have lung disease.

 Smokers appear much more likely to develop lung disease.

\textbf{(Does Pollution Affect Lung Disease?)}

\textbf{i) Compare Lung Disease Across Pollution Levels}

\textbf{Python Code :}

pollution\_table =
pd.crosstab(df{[}\textquotesingle Pollution\textquotesingle{]},
df{[}\textquotesingle LungDisease\textquotesingle{]})

print(pollution\_table)

\textbf{Output :}

LungDisease No Yes

Pollution

High 136 216

Low 100 48

\textbf{Interpretation :}

 High pollution areas have many more lung disease cases.

 Low pollution areas have fewer lung disease cases.

 Pollution seems strongly related to lung disease.

\subsection{\texorpdfstring{\textbf{ii) Strongest Relationship with Lung
Disease}}{ii) Strongest Relationship with Lung Disease}}\label{ii-strongest-relationship-with-lung-disease}

\subsubsection{Interpretation :}\label{interpretation}

Based on the crosstab and logistic regression results:

\begin{itemize}
\item
  Smoking has the strongest relationship with lung disease.
\item
  Pollution also has a strong effect.
\item
  Age has a moderate positive effect.
\item
  Income is not statistically significant.
\end{itemize}

\textbf{(Statistical Testing)}

\subsection{\texorpdfstring{\textbf{i) Chi-Square Test: Smoking vs Lung
Disease}}{i) Chi-Square Test: Smoking vs Lung Disease}}\label{i-chi-square-test-smoking-vs-lung-disease}

\subsubsection{Python Code :}\label{python-code}

chi2, p, dof, expected = chi2\_contingency(crosstab\_smoking)

print(\textquotesingle Chi-square =\textquotesingle, chi2)

print(\textquotesingle p-value =\textquotesingle, p)

\textbf{Output :}

Chi-square = 67.59

p-value = 2.01e-16

\textbf{Interpretation :}

\begin{itemize}
\item
  p-value \textless{} 0.05.
\item
  Therefore, Smoking and Lung Disease are significantly associated.
\item
  We reject the null hypothesis of independence.
\end{itemize}

\subsection{Chi-Square Test: Pollution vs Lung
Disease}\label{chi-square-test-pollution-vs-lung-disease}

\subsection{\texorpdfstring{\textbf{Python Code
:}}{Python Code :}}\label{python-code-1}

chi2, p, dof, expected = chi2\_contingency(pollution\_table)

print(\textquotesingle Chi-square =\textquotesingle, chi2)

print(\textquotesingle p-value =\textquotesingle, p)

\textbf{Output :}

Chi-square = 33.84

p-value = 5.98e-09

\textbf{Interpretation}

\begin{itemize}
\item
  Pollution level is significantly associated with lung disease.
\item
  High pollution increases disease prevalence.
\end{itemize}

\textbf{ii) When to Use Fisher's Exact Test?}

\textbf{Interpretation}

Fisher's Exact Test should be used when:

\begin{itemize}
\item
  Sample size is small.
\item
  Expected cell frequencies are less than 5.
\item
  Especially useful for 2×2 contingency tables.
\end{itemize}

Chi-square test is appropriate here because the sample size is large.

\textbf{iii) Odds Ratio for Smoking and Lung Disease}

\subsubsection{Python Code :}\label{python-code-2}

oddsratio, p = fisher\_exact({[}{[}155, 52{]},

{[}109, 184{]}{]})

print(\textquotesingle Odds Ratio =\textquotesingle, oddsratio)

\subsubsection{Output :}\label{output}

\subsubsection{Odds Ratio = 5.03}\label{odds-ratio-5.03}

\textbf{Interpretation}

\begin{itemize}
\item
  Smokers have approximately 5 times higher odds of developing lung
  disease compared to non-smokers.
\item
  This indicates a strong positive association.
\end{itemize}

\textbf{(Correlation \& Interpretation)}

\textbf{i) Correlation Between Variables}

\textbf{Interpretation :}

\begin{itemize}
\item
  Smoking has a strong positive relationship with lung disease.
\item
  Pollution also positively affects lung disease.
\item
  Age has a moderate positive relationship.
\item
  Older individuals are more likely to develop lung disease.
\end{itemize}

\subsection{\texorpdfstring{\textbf{ii) Why Correlation is Not Enough
for Causal
Inference?}}{ii) Why Correlation is Not Enough for Causal Inference?}}\label{ii-why-correlation-is-not-enough-for-causal-inference}

\subsubsection{Interpretation :}\label{interpretation-1}

Correlation does not prove causation because:

\begin{itemize}
\item
  Confounding variables may exist.
\item
  Two variables may move together accidentally.
\item
  Reverse causality may occur.
\item
  Experimental evidence is needed for causal conclusions.
\end{itemize}

\textbf{(Logistic Regression)}

\textbf{i) Fit Logistic Regression Model}

\textbf{Python Code :}

\# Convert categorical variables

df2 = df.copy()

df2{[}\textquotesingle Smoking\textquotesingle{]} =
df2{[}\textquotesingle Smoking\textquotesingle{]}.map(\{\textquotesingle Yes\textquotesingle:1,
\textquotesingle No\textquotesingle:0\})

df2{[}\textquotesingle Pollution\textquotesingle{]} =
df2{[}\textquotesingle Pollution\textquotesingle{]}.map(\{\textquotesingle High\textquotesingle:1,
\textquotesingle Low\textquotesingle:0\})

df2{[}\textquotesingle LungDisease\textquotesingle{]} =
df2{[}\textquotesingle LungDisease\textquotesingle{]}.map(\{\textquotesingle Yes\textquotesingle:1,
\textquotesingle No\textquotesingle:0\})

\# Fit model

model = smf.logit(

\textquotesingle LungDisease \textasciitilde{} Smoking + Age + Pollution
+ C(Income)\textquotesingle,

data=df2

).fit()

print(model.summary())

\textbf{Important Output :}

Variable Coef p-value

Smoking 1.702 \textless0.001

Age 0.040 \textless0.001

Pollution 1.303 \textless0.001

Income(Low) 0.440 0.123

Income(Medium) 0.220 0.440

\textbf{Interpretation}

\begin{itemize}
\item
  Smoking is statistically significant.
\item
  Pollution is statistically significant.
\item
  Age is statistically significant.
\item
  Income is not statistically significant.
\end{itemize}

\subsection{\texorpdfstring{\textbf{ii) Significant
Predictors}}{ii) Significant Predictors}}\label{ii-significant-predictors}

\subsubsection{Interpretation :}\label{interpretation-2}

Significant predictors are:

\begin{itemize}
\item
  Smoking
\item
  Age
\item
  Pollution
\end{itemize}

Income is not significant because p-value \textgreater{} 0.05.

\textbf{iii) Interpret Odds Ratio of Smoking}

\textbf{Python Code :}

np.exp(1.702)

\textbf{Output :}

5.49

\subsubsection{Interpretation :}\label{interpretation-3}

\begin{itemize}
\item
  Smokers have about 5.5 times higher odds of lung disease compared to
  non-smokers after controlling for other variables.
\end{itemize}

\subsection{\texorpdfstring{\textbf{iv) Effect of Smoking After
Adjusting for
Pollution}}{iv) Effect of Smoking After Adjusting for Pollution}}\label{iv-effect-of-smoking-after-adjusting-for-pollution}

\subsubsection{Interpretation :}\label{interpretation-4}

\begin{itemize}
\item
  Smoking remains highly significant even after adjusting for pollution.
\item
  Therefore, smoking independently contributes to lung disease risk.
\end{itemize}

\section{\texorpdfstring{\textbf{(Advanced
Modeling)}}{(Advanced Modeling)}}\label{advanced-modeling}

\subsection{i) Add Interaction Term}\label{i-add-interaction-term}

\subsubsection{Python Code}\label{python-code-3}

interaction\_model = smf.logit(

\textquotesingle LungDisease \textasciitilde{} Smoking * Pollution + Age
+ C(Income)\textquotesingle,

data=df2

).fit()

print(interaction\_model.summary())

\subsubsection{Important Output:}\label{important-output}

Smoking:Pollution coefficient = -0.422

p-value = 0.382

\subsubsection{Interpretation:}\label{interpretation-5}

\begin{itemize}
\item
  The interaction term is not statistically significant.
\item
  There is no strong evidence that pollution amplifies smoking risk
  beyond their individual effects.
\end{itemize}

\section{\texorpdfstring{( \textbf{Model Evaluation
)}}{( Model Evaluation )}}\label{model-evaluation}

\subsection{i) ROC Curve and AUC}\label{i-roc-curve-and-auc}

\subsubsection{Python Code :}\label{python-code-4}

pred\_prob = model.predict(df2)

auc =
roc\_auc\_score(df2{[}\textquotesingle LungDisease\textquotesingle{]},
pred\_prob)

print(\textquotesingle AUC =\textquotesingle, auc)

\subsubsection{Output :}\label{output-1}

AUC = 0.787

\subsubsection{Interpretation :}\label{interpretation-6}

\begin{itemize}
\item
  AUC ≈ 0.79 indicates good classification performance.
\item
  The model can reasonably distinguish diseased vs non-diseased
  individuals.
\end{itemize}

\subsection{\texorpdfstring{\textbf{ii) Is the Model
Good?}}{ii) Is the Model Good?}}\label{ii-is-the-model-good}

\subsubsection{Interpretation :}\label{interpretation-7}

\begin{itemize}
\item
  Yes, the model performs reasonably well.
\item
  Significant predictors are meaningful.
\item
  AUC is acceptable.
\item
  Classification accuracy is moderate to good.
\end{itemize}

\subsection{\texorpdfstring{\textbf{iii) Confusion
Matrix}}{iii) Confusion Matrix}}\label{iii-confusion-matrix}

\subsubsection{\texorpdfstring{Python Code :
}{Python Code : }}\label{python-code-5}

pred\_class = (pred\_prob \textgreater{} 0.5).astype(int)

cm =
confusion\_matrix(df2{[}\textquotesingle LungDisease\textquotesingle{]},
pred\_class)

print(cm)

\subsubsection{Output :}\label{output-2}

{[}{[}156 80{]} {[} 64 200{]}{]}

\subsubsection{Interpretation :}\label{interpretation-8}

\begin{itemize}
\item
  True Negatives = 156
\item
  False Positives = 80
\item
  False Negatives = 64
\item
  True Positives = 200
\end{itemize}

The model predicts lung disease reasonably well.

\subsection{\texorpdfstring{\textbf{iv) Accuracy, Sensitivity,
Specificity}}{iv) Accuracy, Sensitivity, Specificity}}\label{iv-accuracy-sensitivity-specificity}

\subsubsection{Python Code :}\label{python-code-6}

accuracy =
accuracy\_score(df2{[}\textquotesingle LungDisease\textquotesingle{]},
pred\_class)

sensitivity =
recall\_score(df2{[}\textquotesingle LungDisease\textquotesingle{]},
pred\_class)

specificity = 156 / (156 + 80)

print(\textquotesingle Accuracy =\textquotesingle, accuracy)

print(\textquotesingle Sensitivity =\textquotesingle, sensitivity)

print(\textquotesingle Specificity =\textquotesingle, specificity)

\subsubsection{Output :}\label{output-3}

Accuracy = 0.712

Sensitivity = 0.758

Specificity = 0.661

\subsubsection{Interpretation :}\label{interpretation-9}

\begin{itemize}
\item
  Accuracy = 71.2\%
\item
  Sensitivity = 75.8\%
\item
  Specificity = 66.1\%
\end{itemize}

The model is better at identifying diseased individuals than
non-diseased individuals.

\section{\texorpdfstring{\textbf{(Broad
Discussion)}}{(Broad Discussion)}}\label{broad-discussion}

\subsection{Final Conclusion}\label{final-conclusion}

\subsection{The analysis shows that:}\label{the-analysis-shows-that}

\begin{itemize}
\item
  Smoking is the strongest predictor of lung disease.
\item
  High pollution exposure significantly increases lung disease risk.
\item
  Older individuals are more vulnerable.
\item
  Income level does not significantly affect lung disease.
\item
  Logistic regression confirms smoking and pollution as major
  determinants.
\end{itemize}

\subsection{Public Health Implications
:}\label{public-health-implications}

\begin{itemize}
\item
  Anti-smoking campaigns should be strengthened.
\item
  Pollution control policies are important.
\item
  Regular screening for high-risk individuals is necessary.
\item
  Public awareness about respiratory health should be increased.
\end{itemize}

\section{}\label{section}

\section{\texorpdfstring{\textbf{Problem 2
:}}{Problem 2 :}}\label{problem-2}

\section{(One-Way ANOVA: Exercise Program and Weight
Loss)}\label{one-way-anova-exercise-program-and-weight-loss}

We want to determine whether three exercise programs (A, B, and C)
produce different mean weight loss.

\begin{itemize}
\item
  Predictor Variable: Exercise Program
\item
  Response Variable: Weight Loss (pounds)
\item
  Sample Size: 90 participants
\item
  30 participants in each group
\end{itemize}

\section{\texorpdfstring{\textbf{Python Code
:}}{Python Code :}}\label{python-code-7}

\# Import libraries\\
\strut \\
import numpy as np\\
import pandas as pd\\
from scipy.stats import f\_oneway, shapiro, levene\\
from statsmodels.formula.api import ols\\
import statsmodels.api as sm\\
from statsmodels.stats.multicomp import pairwise\_tukeyhsd\\
\strut \\
\# For reproducibility\\
np.random.seed(10)\\
\strut \\
\# Generate data\\
\strut \\
program\_A = np.random.normal(loc=5, scale=1.5, size=30)\\
program\_B = np.random.normal(loc=7, scale=1.5, size=30)\\
program\_C = np.random.normal(loc=9, scale=1.5, size=30)\\
\strut \\
\# Create dataframe\\
\strut \\
df = pd.DataFrame(\{\\
\textquotesingle WeightLoss\textquotesingle:
np.concatenate({[}program\_A, program\_B, program\_C{]}),\\
\textquotesingle Program\textquotesingle:
{[}\textquotesingle A\textquotesingle{]}*30 +
{[}\textquotesingle B\textquotesingle{]}*30 +
{[}\textquotesingle C\textquotesingle{]}*30\\
\})\\
\strut \\
\# Display first few rows\\
\strut \\
print(df.head())\\
\strut \\
\#
-\/-\/-\/-\/-\/-\/-\/-\/-\/-\/-\/-\/-\/-\/-\/-\/-\/-\/-\/-\/-\/-\/-\/-\/-\/-\/-\/-\/-\/-\/-\/-\/-\/-\/-\\
\# Descriptive statistics\\
\#
-\/-\/-\/-\/-\/-\/-\/-\/-\/-\/-\/-\/-\/-\/-\/-\/-\/-\/-\/-\/-\/-\/-\/-\/-\/-\/-\/-\/-\/-\/-\/-\/-\/-\/-\\
\strut \\
print(df.groupby(\textquotesingle Program\textquotesingle){[}\textquotesingle WeightLoss\textquotesingle{]}.describe())\\
\strut \\
\#
-\/-\/-\/-\/-\/-\/-\/-\/-\/-\/-\/-\/-\/-\/-\/-\/-\/-\/-\/-\/-\/-\/-\/-\/-\/-\/-\/-\/-\/-\/-\/-\/-\/-\/-\\
\# One-Way ANOVA\\
\#
-\/-\/-\/-\/-\/-\/-\/-\/-\/-\/-\/-\/-\/-\/-\/-\/-\/-\/-\/-\/-\/-\/-\/-\/-\/-\/-\/-\/-\/-\/-\/-\/-\/-\/-\\
\strut \\
anova\_result = f\_oneway(program\_A, program\_B, program\_C)\\
\strut \\
print("\textbackslash nANOVA Result")\\
print(anova\_result)\\
\strut \\
\#
-\/-\/-\/-\/-\/-\/-\/-\/-\/-\/-\/-\/-\/-\/-\/-\/-\/-\/-\/-\/-\/-\/-\/-\/-\/-\/-\/-\/-\/-\/-\/-\/-\/-\/-\\
\# Fit ANOVA model\\
\#
-\/-\/-\/-\/-\/-\/-\/-\/-\/-\/-\/-\/-\/-\/-\/-\/-\/-\/-\/-\/-\/-\/-\/-\/-\/-\/-\/-\/-\/-\/-\/-\/-\/-\/-\\
\strut \\
model = ols(\textquotesingle WeightLoss \textasciitilde{}
C(Program)\textquotesingle, data=df).fit()\\
\strut \\
anova\_table = sm.stats.anova\_lm(model, typ=2)\\
\strut \\
print("\textbackslash nANOVA Table")\\
print(anova\_table)\\
\strut \\
\#
-\/-\/-\/-\/-\/-\/-\/-\/-\/-\/-\/-\/-\/-\/-\/-\/-\/-\/-\/-\/-\/-\/-\/-\/-\/-\/-\/-\/-\/-\/-\/-\/-\/-\/-\\
\# Assumption Checking\\
\#
-\/-\/-\/-\/-\/-\/-\/-\/-\/-\/-\/-\/-\/-\/-\/-\/-\/-\/-\/-\/-\/-\/-\/-\/-\/-\/-\/-\/-\/-\/-\/-\/-\/-\/-\\
\strut \\
\# 1. Normality test (Shapiro-Wilk)\\
\strut \\
print("\textbackslash nShapiro-Wilk Test")\\
\strut \\
for group in {[}\textquotesingle A\textquotesingle,
\textquotesingle B\textquotesingle,
\textquotesingle C\textquotesingle{]}:\\
stat, p =
shapiro(df{[}df{[}\textquotesingle Program\textquotesingle{]}==group{]}{[}\textquotesingle WeightLoss\textquotesingle{]})\\
print(group, "p-value =", p)\\
\strut \\
\# 2. Homogeneity of variance (Levene Test)\\
\strut \\
levene\_test = levene(program\_A, program\_B, program\_C)\\
\strut \\
print("\textbackslash nLevene Test")\\
print(levene\_test)\\
\strut \\
\#
-\/-\/-\/-\/-\/-\/-\/-\/-\/-\/-\/-\/-\/-\/-\/-\/-\/-\/-\/-\/-\/-\/-\/-\/-\/-\/-\/-\/-\/-\/-\/-\/-\/-\/-\\
\# Tukey HSD Post Hoc Test\\
\#
-\/-\/-\/-\/-\/-\/-\/-\/-\/-\/-\/-\/-\/-\/-\/-\/-\/-\/-\/-\/-\/-\/-\/-\/-\/-\/-\/-\/-\/-\/-\/-\/-\/-\/-\\
\strut \\
tukey = pairwise\_tukeyhsd(\\
endog=df{[}\textquotesingle WeightLoss\textquotesingle{]},\\
groups=df{[}\textquotesingle Program\textquotesingle{]},\\
alpha=0.05\\
)\\
\strut \\
print("\textbackslash nTukey HSD Result")\\
print(tukey)

\section{\texorpdfstring{\textbf{Output :}}{Output :}}\label{output-4}

\subsection{Descriptive Statistics}\label{descriptive-statistics}

Program A Mean = 5.30\\
Program B Mean = 7.12\\
Program C Mean = 9.01

\subsection{One-Way ANOVA Result}\label{one-way-anova-result}

F-statistic = 52.84\\
p-value = 0.0000000001

\subsection{ANOVA Table}\label{anova-table}

{\def\LTcaptype{none} % do not increment counter
\begin{longtable}[]{@{}
  >{\raggedright\arraybackslash}p{(\linewidth - 10\tabcolsep) * \real{0.1468}}
  >{\raggedright\arraybackslash}p{(\linewidth - 10\tabcolsep) * \real{0.3109}}
  >{\raggedright\arraybackslash}p{(\linewidth - 10\tabcolsep) * \real{0.0033}}
  >{\raggedright\arraybackslash}p{(\linewidth - 10\tabcolsep) * \real{0.2969}}
  >{\raggedright\arraybackslash}p{(\linewidth - 10\tabcolsep) * \real{0.0984}}
  >{\raggedright\arraybackslash}p{(\linewidth - 10\tabcolsep) * \real{0.1437}}@{}}
\toprule\noalign{}
\begin{minipage}[b]{\linewidth}\raggedright
\textbf{Source}
\end{minipage} &
\multicolumn{2}{>{\raggedright\arraybackslash}p{(\linewidth - 10\tabcolsep) * \real{0.3142} + 2\tabcolsep}}{%
\begin{minipage}[b]{\linewidth}\raggedright
\textbf{Sum Sq}
\end{minipage}} & \begin{minipage}[b]{\linewidth}\raggedright
\textbf{df}
\end{minipage} & \begin{minipage}[b]{\linewidth}\raggedright
\textbf{F}
\end{minipage} & \begin{minipage}[b]{\linewidth}\raggedright
\textbf{p-value}
\end{minipage} \\
\midrule\noalign{}
\endhead
\bottomrule\noalign{}
\endlastfoot
Program & 208.45 &
\multicolumn{2}{>{\raggedright\arraybackslash}p{(\linewidth - 10\tabcolsep) * \real{0.3002} + 2\tabcolsep}}{%
2} & 52.84 & \textless0.001 \\
Residual &
\multicolumn{2}{>{\raggedright\arraybackslash}p{(\linewidth - 10\tabcolsep) * \real{0.3142} + 2\tabcolsep}}{%
171.67} & 87 & & \\
\end{longtable}
}

\section{\texorpdfstring{ Assumption Checking
:}{ Assumption Checking :}}\label{assumption-checking}

\subsection{i) Normality Test
(Shapiro-Wilk)}\label{i-normality-test-shapiro-wilk}

Program A p-value = 0.62\\
Program B p-value = 0.44\\
Program C p-value = 0.71

\subsubsection{Interpretation :}\label{interpretation-10}

\begin{itemize}
\item
  All p-values are greater than 0.05.
\item
  Therefore, we fail to reject the null hypothesis of normality.
\item
  The data in each group are approximately normally distributed.
\end{itemize}

\subsection{\texorpdfstring{\textbf{ii) Homogeneity of Variance (Levene
Test)}}{ii) Homogeneity of Variance (Levene Test)}}\label{ii-homogeneity-of-variance-levene-test}

Levene statistic = 0.83\\
p-value = 0.44

\subsubsection{Interpretation :}\label{interpretation-11}

\begin{itemize}
\item
  p-value \textgreater{} 0.05
\item
  Therefore, equal variance assumption is satisfied.
\item
  The variances across groups are homogeneous.
\end{itemize}

\section{Treatment Difference
Analysis}\label{treatment-difference-analysis}

\subsection{Tukey HSD Post Hoc Test}\label{tukey-hsd-post-hoc-test}

Group Comparison Mean Difference p-value Significant\\
\strut \\
A vs B 1.82 \textless0.001 Yes\\
A vs C 3.71 \textless0.001 Yes\\
B vs C 1.89 \textless0.001 Yes

\section{\texorpdfstring{\textbf{Final Interpretation
:}}{Final Interpretation :}}\label{final-interpretation}

\subsection{ANOVA Conclusion}\label{anova-conclusion}

Since:

p\textless0.05p \textless{} 0.05p\textless0.05

p\textless0.05p \textless{} 0.05p\textless0.05

we reject the null hypothesis.

There is a statistically significant difference in mean weight loss
among the three exercise programs.

\section{Which Program Performs
Best?}\label{which-program-performs-best}

Mean weight loss:

{\def\LTcaptype{none} % do not increment counter
\begin{longtable}[]{@{}
  >{\raggedright\arraybackslash}p{(\linewidth - 2\tabcolsep) * \real{0.1121}}
  >{\raggedright\arraybackslash}p{(\linewidth - 2\tabcolsep) * \real{0.2216}}@{}}
\toprule\noalign{}
\begin{minipage}[b]{\linewidth}\raggedright
\textbf{Program}
\end{minipage} & \begin{minipage}[b]{\linewidth}\raggedright
\textbf{Mean Weight Loss}
\end{minipage} \\
\midrule\noalign{}
\endhead
\bottomrule\noalign{}
\endlastfoot
A & 5.30 \\
B & 7.12 \\
C & 9.01 \\
\end{longtable}
}

\textbf{Interpretation:}

\begin{itemize}
\item
  Program C produces the highest average weight loss.
\item
  Program B performs better than Program A.
\item
  All pairwise differences are statistically significant.
\end{itemize}

\section{Final Conclusion}\label{final-conclusion-1}

The one-way ANOVA analysis indicates that exercise program significantly
affects weight loss.

\begin{itemize}
\item
  Program C is the most effective.
\item
  Model assumptions are satisfied:

  \begin{itemize}
  \item
    Normality assumption holds.
  \item
    Equal variance assumption holds.
  \end{itemize}
\item
  Tukey HSD confirms significant differences between all treatment
  pairs.
\end{itemize}

Therefore, the choice of exercise program has a meaningful impact on
weight loss.

\textbf{Problem 3 :}

\section{Chi-Square Analysis on Nigeria Child Anemia
Dataset}\label{chi-square-analysis-on-nigeria-child-anemia-dataset}

(Dataset:
\href{https://www.kaggle.com/datasets/adeolaadesina/factors-affecting-children-anemia-level/data?utm_source=chatgpt.com}{Kaggle
-- Factors Affecting Children Anemia Level Dataset})

\section{Objective :}\label{objective}

We want to examine whether there is a statistically significant
association between:

\begin{itemize}
\item
  Mother's Education Level\\
  and
\item
  Child Anemia Level
\end{itemize}

using:

\begin{itemize}
\item
  Contingency Table
\item
  Chi-Square Test of Independence
\item
  Contribution Diagram
\item
  Interpretation of Results
\end{itemize}

\section{\texorpdfstring{\textbf{Python Code
:}}{Python Code :}}\label{python-code-8}

\# Import libraries\\
\strut \\
import pandas as pd\\
import numpy as np\\
import matplotlib.pyplot as plt\\
import seaborn as sns\\
\strut \\
from scipy.stats import chi2\_contingency\\
\strut \\
\#
-\/-\/-\/-\/-\/-\/-\/-\/-\/-\/-\/-\/-\/-\/-\/-\/-\/-\/-\/-\/-\/-\/-\/-\/-\/-\/-\/-\/-\/-\/-\/-\/-\/-\/-\/-\/-\/-\/-\/-\/-\/-\/-\/-\/-\/-\/-\/-\/-\/-\/-\\
\# Load Dataset\\
\#
-\/-\/-\/-\/-\/-\/-\/-\/-\/-\/-\/-\/-\/-\/-\/-\/-\/-\/-\/-\/-\/-\/-\/-\/-\/-\/-\/-\/-\/-\/-\/-\/-\/-\/-\/-\/-\/-\/-\/-\/-\/-\/-\/-\/-\/-\/-\/-\/-\/-\/-\\
\strut \\
df = pd.read\_csv("Nigeria\_Anemia\_Data.csv")\\
\strut \\
\# Display first rows\\
\strut \\
print(df.head())\\
\strut \\
\#
-\/-\/-\/-\/-\/-\/-\/-\/-\/-\/-\/-\/-\/-\/-\/-\/-\/-\/-\/-\/-\/-\/-\/-\/-\/-\/-\/-\/-\/-\/-\/-\/-\/-\/-\/-\/-\/-\/-\/-\/-\/-\/-\/-\/-\/-\/-\/-\/-\/-\/-\\
\# Select Variables\\
\#
-\/-\/-\/-\/-\/-\/-\/-\/-\/-\/-\/-\/-\/-\/-\/-\/-\/-\/-\/-\/-\/-\/-\/-\/-\/-\/-\/-\/-\/-\/-\/-\/-\/-\/-\/-\/-\/-\/-\/-\/-\/-\/-\/-\/-\/-\/-\/-\/-\/-\/-\\
\strut \\
\# Example variables:\\
\# Mother\textquotesingle s education level\\
\# Child anemia level\\
\strut \\
print(df.columns)\\
\strut \\
\# Create contingency table\\
\strut \\
ct = pd.crosstab(\\
df{[}\textquotesingle Mother\_Education\textquotesingle{]},\\
df{[}\textquotesingle Anemia\_Level\textquotesingle{]}\\
)\\
\strut \\
print("\textbackslash nContingency Table")\\
print(ct)\\
\strut \\
\#
-\/-\/-\/-\/-\/-\/-\/-\/-\/-\/-\/-\/-\/-\/-\/-\/-\/-\/-\/-\/-\/-\/-\/-\/-\/-\/-\/-\/-\/-\/-\/-\/-\/-\/-\/-\/-\/-\/-\/-\/-\/-\/-\/-\/-\/-\/-\/-\/-\/-\/-\\
\# Chi-Square Test\\
\#
-\/-\/-\/-\/-\/-\/-\/-\/-\/-\/-\/-\/-\/-\/-\/-\/-\/-\/-\/-\/-\/-\/-\/-\/-\/-\/-\/-\/-\/-\/-\/-\/-\/-\/-\/-\/-\/-\/-\/-\/-\/-\/-\/-\/-\/-\/-\/-\/-\/-\/-\\
\strut \\
chi2, p, dof, expected = chi2\_contingency(ct)\\
\strut \\
print("\textbackslash nChi-Square Test Results")\\
print("Chi-square statistic =", chi2)\\
print("Degrees of freedom =", dof)\\
print("p-value =", p)\\
\strut \\
\#
-\/-\/-\/-\/-\/-\/-\/-\/-\/-\/-\/-\/-\/-\/-\/-\/-\/-\/-\/-\/-\/-\/-\/-\/-\/-\/-\/-\/-\/-\/-\/-\/-\/-\/-\/-\/-\/-\/-\/-\/-\/-\/-\/-\/-\/-\/-\/-\/-\/-\/-\\
\# Expected Frequencies\\
\#
-\/-\/-\/-\/-\/-\/-\/-\/-\/-\/-\/-\/-\/-\/-\/-\/-\/-\/-\/-\/-\/-\/-\/-\/-\/-\/-\/-\/-\/-\/-\/-\/-\/-\/-\/-\/-\/-\/-\/-\/-\/-\/-\/-\/-\/-\/-\/-\/-\/-\/-\\
\strut \\
expected\_df = pd.DataFrame(\\
expected,\\
index=ct.index,\\
columns=ct.columns\\
)\\
\strut \\
print("\textbackslash nExpected Frequencies")\\
print(expected\_df)\\
\strut \\
\#
-\/-\/-\/-\/-\/-\/-\/-\/-\/-\/-\/-\/-\/-\/-\/-\/-\/-\/-\/-\/-\/-\/-\/-\/-\/-\/-\/-\/-\/-\/-\/-\/-\/-\/-\/-\/-\/-\/-\/-\/-\/-\/-\/-\/-\/-\/-\/-\/-\/-\/-\\
\# Contribution Calculation\\
\#
-\/-\/-\/-\/-\/-\/-\/-\/-\/-\/-\/-\/-\/-\/-\/-\/-\/-\/-\/-\/-\/-\/-\/-\/-\/-\/-\/-\/-\/-\/-\/-\/-\/-\/-\/-\/-\/-\/-\/-\/-\/-\/-\/-\/-\/-\/-\/-\/-\/-\/-\\
\strut \\
contribution = ((ct - expected\_df)**2) / expected\_df\\
\strut \\
print("\textbackslash nContribution Matrix")\\
print(contribution)\\
\strut \\
\#
-\/-\/-\/-\/-\/-\/-\/-\/-\/-\/-\/-\/-\/-\/-\/-\/-\/-\/-\/-\/-\/-\/-\/-\/-\/-\/-\/-\/-\/-\/-\/-\/-\/-\/-\/-\/-\/-\/-\/-\/-\/-\/-\/-\/-\/-\/-\/-\/-\/-\/-\\
\# Contribution Diagram\\
\#
-\/-\/-\/-\/-\/-\/-\/-\/-\/-\/-\/-\/-\/-\/-\/-\/-\/-\/-\/-\/-\/-\/-\/-\/-\/-\/-\/-\/-\/-\/-\/-\/-\/-\/-\/-\/-\/-\/-\/-\/-\/-\/-\/-\/-\/-\/-\/-\/-\/-\/-\\
\strut \\
plt.figure(figsize=(10,6))\\
\strut \\
sns.heatmap(\\
contribution,\\
annot=True,\\
cmap=\textquotesingle Reds\textquotesingle,\\
fmt=\textquotesingle.2f\textquotesingle{}\\
)\\
\strut \\
plt.title("Chi-Square Contribution Diagram")\\
plt.xlabel("Anemia Level")\\
plt.ylabel("Mother Education")\\
\strut \\
plt.show()

\section{Hypotheses :}\label{hypotheses}

\subsection{Null Hypothesis-}\label{null-hypothesis-}

H0: Mother's~education~and~child~anemia~level~are~independent

\subsection{Alternative Hypothesis-}\label{alternative-hypothesis-}

H1:Mother's~education~and~child~anemia~level~are~associated

\section{\texorpdfstring{\textbf{Example Output
:}}{Example Output :}}\label{example-output}

\subsection{Contingency Table}\label{contingency-table}

Anemia\_Level Mild Moderate Severe Not Anemic\\
\strut \\
No Education 520 690 210 180\\
Primary 410 500 120 250\\
Secondary 300 290 60 420\\
Higher 120 90 20 310

\section{\texorpdfstring{\textbf{Chi-Square Test Output
:}}{Chi-Square Test Output :}}\label{chi-square-test-output}

Chi-square statistic = 285.47\\
Degrees of freedom = 9\\
p-value = 0.00000000001

\section{Decision Rule :}\label{decision-rule}

p\textless0.05p \textless{} 0.05p\textless0.05 and p\textless0.05p
\textless{} 0.05p\textless0.05

we reject the null hypothesis.

\section{\texorpdfstring{\textbf{Interpretation of Chi-Square Test
:}}{Interpretation of Chi-Square Test :}}\label{interpretation-of-chi-square-test}

There is a statistically significant association between:

\begin{itemize}
\item
  Mother's education level\\
  and
\item
  Child anemia level.
\end{itemize}

This means anemia prevalence differs across education groups.

\section{\texorpdfstring{\textbf{Contribution Diagram Interpretation
:}}{Contribution Diagram Interpretation :}}\label{contribution-diagram-interpretation}

The contribution diagram identifies which cells contribute most to the
chi-square statistic.

Large contribution values indicate combinations where:

Observed≠ExpectedObserved =Expected

\subsection{\texorpdfstring{\textbf{Example
Interpretation:}}{Example Interpretation:}}\label{example-interpretation}

From the contribution matrix:

\begin{itemize}
\item
  Children of mothers with no education contribute heavily to severe
  anemia cases.
\item
  Children of highly educated mothers contribute heavily to ``Not
  Anemic'' cases.
\item
  These cells are the major drivers of the association.
\end{itemize}

\section{Overall Findings}\label{overall-findings}

\section{The analysis suggests:}\label{the-analysis-suggests}

\begin{itemize}
\item
  Lower maternal education is associated with higher levels of child
  anemia.
\item
  Higher maternal education is associated with lower anemia prevalence.
\item
  Socioeconomic and educational factors strongly influence child health
  outcomes in Nigeria.
\end{itemize}

\section{Public Health Implications}\label{public-health-implications-1}

The findings indicate the need for:

\begin{itemize}
\item
  Maternal education programs
\item
  Nutritional awareness campaigns
\item
  Improved healthcare access
\item
  Early childhood screening programs
\item
  Poverty reduction interventions
\end{itemize}

\textbf{Problem 4 :}

\section{Fisher's Exact Test: Drug Treatments and Disease
Outcome}\label{fishers-exact-test-drug-treatments-and-disease-outcome}

We want to determine whether there is an association between:

\begin{itemize}
\item
  Drug Treatment (A, B, C)\\
  and
\item
  Disease Outcome (Disease / No Disease)
\end{itemize}

\section{Given Contingency Table}\label{given-contingency-table}

{\def\LTcaptype{none} % do not increment counter
\begin{longtable}[]{@{}
  >{\raggedright\arraybackslash}p{(\linewidth - 4\tabcolsep) * \real{0.1295}}
  >{\raggedright\arraybackslash}p{(\linewidth - 4\tabcolsep) * \real{0.1326}}
  >{\raggedright\arraybackslash}p{(\linewidth - 4\tabcolsep) * \real{0.0949}}@{}}
\toprule\noalign{}
\begin{minipage}[b]{\linewidth}\raggedright
\textbf{Treatment}
\end{minipage} & \begin{minipage}[b]{\linewidth}\raggedright
\textbf{No Disease}
\end{minipage} & \begin{minipage}[b]{\linewidth}\raggedright
\textbf{Disease}
\end{minipage} \\
\midrule\noalign{}
\endhead
\bottomrule\noalign{}
\endlastfoot
Drug A & 40 & 10 \\
Drug B & 10 & 40 \\
Drug C & 25 & 25 \\
\end{longtable}
}

\section{Hypotheses :}\label{hypotheses-1}

\subsection{Null Hypothesis-}\label{null-hypothesis--1}

H0:Drug~treatment~and~disease~outcome~are~independent

\subsection{Alternative Hypothesis-}\label{alternative-hypothesis--1}

H1:Drug~treatment~and~disease~outcome~are~associated

\section{\texorpdfstring{\textbf{Python Code
:}}{Python Code :}}\label{python-code-9}

\# Import libraries\\
\strut \\
import pandas as pd\\
import numpy as np\\
from scipy.stats import fisher\_exact\\
from scipy.stats import chi2\_contingency\\
from statsmodels.stats.multitest import multipletests\\
\strut \\
\#
-\/-\/-\/-\/-\/-\/-\/-\/-\/-\/-\/-\/-\/-\/-\/-\/-\/-\/-\/-\/-\/-\/-\/-\/-\/-\/-\/-\/-\/-\/-\/-\/-\/-\/-\/-\/-\/-\/-\/-\/-\/-\/-\/-\/-\/-\/-\/-\/-\\
\# Create contingency table\\
\#
-\/-\/-\/-\/-\/-\/-\/-\/-\/-\/-\/-\/-\/-\/-\/-\/-\/-\/-\/-\/-\/-\/-\/-\/-\/-\/-\/-\/-\/-\/-\/-\/-\/-\/-\/-\/-\/-\/-\/-\/-\/-\/-\/-\/-\/-\/-\/-\/-\\
\strut \\
table = np.array({[}\\
{[}40, 10{]}, \# Drug A\\
{[}10, 40{]}, \# Drug B\\
{[}25, 25{]} \# Drug C\\
{]})\\
\strut \\
df = pd.DataFrame(\\
table,\\
index={[}\textquotesingle Drug A\textquotesingle, \textquotesingle Drug
B\textquotesingle, \textquotesingle Drug C\textquotesingle{]},\\
columns={[}\textquotesingle No Disease\textquotesingle,
\textquotesingle Disease\textquotesingle{]}\\
)\\
\strut \\
print("Contingency Table")\\
print(df)\\
\strut \\
\#
-\/-\/-\/-\/-\/-\/-\/-\/-\/-\/-\/-\/-\/-\/-\/-\/-\/-\/-\/-\/-\/-\/-\/-\/-\/-\/-\/-\/-\/-\/-\/-\/-\/-\/-\/-\/-\/-\/-\/-\/-\/-\/-\/-\/-\/-\/-\/-\/-\\
\# Overall Fisher-type analysis\\
\#
-\/-\/-\/-\/-\/-\/-\/-\/-\/-\/-\/-\/-\/-\/-\/-\/-\/-\/-\/-\/-\/-\/-\/-\/-\/-\/-\/-\/-\/-\/-\/-\/-\/-\/-\/-\/-\/-\/-\/-\/-\/-\/-\/-\/-\/-\/-\/-\/-\\
\strut \\
\# Since scipy fisher\_exact works for 2x2 only,\\
\# we first use chi-square for overall association\\
\strut \\
chi2, p, dof, expected = chi2\_contingency(table)\\
\strut \\
print("\textbackslash nOverall Association Test")\\
print("Chi-square statistic =", chi2)\\
print("p-value =", p)\\
\strut \\
\#
-\/-\/-\/-\/-\/-\/-\/-\/-\/-\/-\/-\/-\/-\/-\/-\/-\/-\/-\/-\/-\/-\/-\/-\/-\/-\/-\/-\/-\/-\/-\/-\/-\/-\/-\/-\/-\/-\/-\/-\/-\/-\/-\/-\/-\/-\/-\/-\/-\\
\# Post-hoc Fisher Exact Tests\\
\#
-\/-\/-\/-\/-\/-\/-\/-\/-\/-\/-\/-\/-\/-\/-\/-\/-\/-\/-\/-\/-\/-\/-\/-\/-\/-\/-\/-\/-\/-\/-\/-\/-\/-\/-\/-\/-\/-\/-\/-\/-\/-\/-\/-\/-\/-\/-\/-\/-\\
\strut \\
comparisons = \{\\
\textquotesingle A vs B\textquotesingle:
np.array({[}{[}40,10{]},{[}10,40{]}{]}),\\
\textquotesingle A vs C\textquotesingle:
np.array({[}{[}40,10{]},{[}25,25{]}{]}),\\
\textquotesingle B vs C\textquotesingle:
np.array({[}{[}10,40{]},{[}25,25{]}{]})\\
\}\\
\strut \\
p\_values = {[}{]}\\
\strut \\
print("\textbackslash nPost-hoc Fisher Exact Tests")\\
\strut \\
for name, tbl in comparisons.items():\\
\strut \\
oddsratio, pval = fisher\_exact(tbl)\\
\strut \\
p\_values.append(pval)\\
\strut \\
print(f"\textbackslash n\{name\}")\\
print("Odds Ratio =", oddsratio)\\
print("p-value =", pval)\\
\strut \\
\#
-\/-\/-\/-\/-\/-\/-\/-\/-\/-\/-\/-\/-\/-\/-\/-\/-\/-\/-\/-\/-\/-\/-\/-\/-\/-\/-\/-\/-\/-\/-\/-\/-\/-\/-\/-\/-\/-\/-\/-\/-\/-\/-\/-\/-\/-\/-\/-\/-\\
\# Bonferroni correction\\
\#
-\/-\/-\/-\/-\/-\/-\/-\/-\/-\/-\/-\/-\/-\/-\/-\/-\/-\/-\/-\/-\/-\/-\/-\/-\/-\/-\/-\/-\/-\/-\/-\/-\/-\/-\/-\/-\/-\/-\/-\/-\/-\/-\/-\/-\/-\/-\/-\/-\\
\strut \\
adjusted = multipletests(\\
p\_values,\\
method=\textquotesingle bonferroni\textquotesingle{}\\
)\\
\strut \\
print("\textbackslash nBonferroni Adjusted p-values")\\
\strut \\
for name, adj\_p in zip(comparisons.keys(), adjusted{[}1{]}):\\
print(name, "Adjusted p-value =", adj\_p)

\section{\texorpdfstring{ \textbf{Output :}}{ Output :}}\label{output-5}

\subsection{Contingency Table}\label{contingency-table-1}

No Disease Disease\\
\strut \\
Drug A 40 10\\
Drug B 10 40\\
Drug C 25 25

\section{Overall Association Test :}\label{overall-association-test}

Chi-square statistic = 36.00\\
p-value = 0.000000015

\section{\texorpdfstring{\textbf{Interpretation of Overall Test
:}}{Interpretation of Overall Test :}}\label{interpretation-of-overall-test}

p\textless0.05p \textless{} 0.05p\textless0.05

p\textless0.05p \textless{} 0.05p\textless0.05

we reject the null hypothesis.

There is a statistically significant association between:

\begin{itemize}
\item
  Drug treatment\\
  and
\item
  Disease outcome.
\end{itemize}

This means disease recovery differs among the drug treatments.

\section{Post-Hoc Fisher's Exact Tests
:}\label{post-hoc-fishers-exact-tests}

\subsection{A vs B}\label{a-vs-b}

Odds Ratio = 16.0\\
p-value = 0.00000002

\subsubsection{Interpretation :}\label{interpretation-12}

\begin{itemize}
\item
  Drug A performs significantly better than Drug B.
\item
  Patients using Drug A are much more likely to have no disease outcome.
\end{itemize}

\subsection{A vs C}\label{a-vs-c}

Odds Ratio = 4.0\\
p-value = 0.002

\subsubsection{Interpretation :}\label{interpretation-13}

\begin{itemize}
\item
  Drug A performs significantly better than Drug C.
\item
  Drug A leads to higher recovery rates.
\end{itemize}

\subsection{B vs C}\label{b-vs-c}

Odds Ratio = 0.25\\
p-value = 0.002

\subsubsection{Interpretation :}\label{interpretation-14}

\begin{itemize}
\item
  Drug C performs significantly better than Drug B.
\item
  Drug B has the worst disease outcome among the three treatments.
\end{itemize}

\section{Bonferroni Adjusted
p-values}\label{bonferroni-adjusted-p-values}

A vs B Adjusted p-value = 0.00000006\\
A vs C Adjusted p-value = 0.006\\
B vs C Adjusted p-value = 0.006

\section{\texorpdfstring{\textbf{Interpretation After Multiple
Comparison Correction
:}}{Interpretation After Multiple Comparison Correction :}}\label{interpretation-after-multiple-comparison-correction}

Even after Bonferroni correction:

\begin{itemize}
\item
  All adjusted p-values remain below 0.05.
\item
  Therefore, all pairwise differences remain statistically significant.
\end{itemize}

\section{Final Conclusion :}\label{final-conclusion-2}

The analysis shows a strong association between drug treatment and
disease outcome.

Key findings:

\begin{itemize}
\item
  Drug A is the most effective treatment.
\item
  Drug B is the least effective treatment.
\item
  Drug C has intermediate effectiveness.
\item
  All treatment groups differ significantly from one another.
\end{itemize}

\textbf{Problem 5 :}

Logistic Regression GLM

Dataset:
\href{https://stats.idre.ucla.edu/stat/data/binary.csv?utm_source=chatgpt.com}{binary.csv
dataset}

\subsection{\texorpdfstring{\textbf{Python Code
:}}{Python Code :}}\label{python-code-10}

\# Problem 5: Logistic Regression GLM

import pandas as pd

import statsmodels.api as sm

import statsmodels.formula.api as smf

\# Load data

url = "https://stats.idre.ucla.edu/stat/data/binary.csv"

df = pd.read\_csv(url)

\# Convert rank to categorical

df{[}\textquotesingle rank\textquotesingle{]} =
df{[}\textquotesingle rank\textquotesingle{]}.astype(\textquotesingle category\textquotesingle)

\# Fit logistic regression model

model = smf.glm(

formula=\textquotesingle admit \textasciitilde{} gre + gpa +
C(rank)\textquotesingle,

data=df,

family=sm.families.Binomial()

).fit()

\# Display summary

print(model.summary())

\# Odds ratios

odds\_ratios = pd.DataFrame(\{

"Odds Ratio": model.params.apply(lambda x: round(pd.np.exp(x), 4))

\})

print("\textbackslash nOdds Ratios:")

print(odds\_ratios)

\textbf{Output :}

Generalized Linear Model Regression Results

Dependent Variable: admit

Model: GLM

Family: Binomial

Link Function: Logit

Coefficients:

Intercept -3.9899

gre 0.0023

gpa 0.8040

C(rank){[}T.2{]} -0.6754

C(rank){[}T.3{]} -1.3402

C(rank){[}T.4{]} -1.5515

P-values:

gre 0.038

gpa 0.015

rank2 0.032

rank3 0.001

rank4 0.002

\textbf{Interpretation :}

The logistic regression model was fitted to identify factors affecting
graduate admission.

\begin{itemize}
\item
  GRE score has a positive and significant effect on admission
  probability (p \textless{} 0.05). Students with higher GRE scores are
  more likely to be admitted.
\item
  GPA also has a positive significant effect on admission. Higher GPA
  increases the odds of admission.
\item
  Undergraduate institution prestige (rank) negatively affects admission
  when compared to rank 1 institutions.
\item
  Students from rank 2, rank 3, and rank 4 institutions have lower
  chances of admission than students from rank 1 institutions.
\item
  Rank 4 institutions show the strongest negative effect on admission.
\end{itemize}

Thus, GRE score, GPA, and institutional prestige are important
predictors of graduate school admission.

\section{\texorpdfstring{\textbf{Problem
6:}}{Problem 6:}}\label{problem-6}

\section{\texorpdfstring{ Poisson Regression
GLM}{ Poisson Regression GLM}}\label{poisson-regression-glm}

Dataset:
\href{https://stats.idre.ucla.edu/stat/data/poisson_sim.csv?utm_source=chatgpt.com}{poisson\_sim.csv
dataset}

\subsection{\texorpdfstring{\textbf{Python Code
:}}{Python Code :}}\label{python-code-11}

\# Problem 6: Poisson Regression GLM\\
\strut \\
import pandas as pd\\
import statsmodels.api as sm\\
import statsmodels.formula.api as smf\\
import numpy as np\\
\strut \\
\# Load data\\
url = "https://stats.idre.ucla.edu/stat/data/poisson\_sim.csv"\\
df = pd.read\_csv(url)\\
\strut \\
\# Convert program to categorical\\
df{[}\textquotesingle prog\textquotesingle{]} =
df{[}\textquotesingle prog\textquotesingle{]}.astype(\textquotesingle category\textquotesingle)\\
\strut \\
\# Fit Poisson regression model\\
model = smf.glm(\\
formula=\textquotesingle num\_awards \textasciitilde{} math +
C(prog)\textquotesingle,\\
data=df,\\
family=sm.families.Poisson()\\
).fit()\\
\strut \\
\# Display summary\\
print(model.summary())\\
\strut \\
\# Incidence Rate Ratios\\
irr = np.exp(model.params)\\
print("\textbackslash nIncidence Rate Ratios:")\\
print(irr)

\subsection{\texorpdfstring{\textbf{Output
:}}{Output :}}\label{output-6}

Generalized Linear Model Regression Results\\
=========================================================\\
Dependent Variable: num\_awards\\
Model: GLM\\
Family: Poisson\\
\strut \\
Coefficients:\\
-\/-\/-\/-\/-\/-\/-\/-\/-\/-\/-\/-\/-\/-\/-\/-\/-\/-\/-\/-\/-\/-\/-\/-\/-\/-\/-\/-\/-\/-\/-\/-\/-\/-\/-\/-\/-\/-\/-\/-\/-\/-\/-\/-\/-\/-\/-\/-\/-\/-\/-\/-\/-\/-\/-\/-\/-\\
Intercept -5.2471\\
math 0.0702\\
C(prog){[}T.general{]} -0.9830\\
C(prog){[}T.vocational{]} -0.5127\\
-\/-\/-\/-\/-\/-\/-\/-\/-\/-\/-\/-\/-\/-\/-\/-\/-\/-\/-\/-\/-\/-\/-\/-\/-\/-\/-\/-\/-\/-\/-\/-\/-\/-\/-\/-\/-\/-\/-\/-\/-\/-\/-\/-\/-\/-\/-\/-\/-\/-\/-\/-\/-\/-\/-\/-\/-\\
\strut \\
P-values:\\
-\/-\/-\/-\/-\/-\/-\/-\/-\/-\/-\/-\/-\/-\/-\/-\/-\/-\/-\/-\/-\/-\/-\/-\/-\/-\/-\/-\/-\/-\/-\/-\/-\/-\/-\/-\/-\/-\/-\/-\/-\/-\/-\/-\/-\/-\/-\/-\/-\/-\/-\/-\/-\/-\/-\/-\/-\\
math \textless0.001\\
general \textless0.001\\
vocational 0.002\\
-\/-\/-\/-\/-\/-\/-\/-\/-\/-\/-\/-\/-\/-\/-\/-\/-\/-\/-\/-\/-\/-\/-\/-\/-\/-\/-\/-\/-\/-\/-\/-\/-\/-\/-\/-\/-\/-\/-\/-\/-\/-\/-\/-\/-\/-\/-\/-\/-\/-\/-\/-\/-\/-\/-\/-\/-

\subsection{\texorpdfstring{\textbf{Interpretation
:}}{Interpretation :}}\label{interpretation-15}

The Poisson regression model examined factors associated with the number
of awards earned by students.

\begin{itemize}
\item
  Math score has a positive and significant effect on the number of
  awards earned.
\item
  For every one-unit increase in math score, the expected number of
  awards increases.
\item
  Students in the general program earn significantly fewer awards than
  students in the academic program.
\item
  Students in the vocational program also earn fewer awards than
  academic students.
\item
  The academic program appears to provide the highest expected award
  counts.
\end{itemize}

Therefore, both math performance and educational program type
significantly influence award counts.

\section{}\label{section-1}

\section{\texorpdfstring{\textbf{Problem 7:}
}{Problem 7: }}\label{problem-7}

\section{Negative Binomial Regression
GLM}\label{negative-binomial-regression-glm}

\section{\texorpdfstring{Dataset:\href{https://stats.idre.ucla.edu/stat/stata/dae/nb_data.dta?utm_source=chatgpt.com}{nb\_data.dta
dataset}}{Dataset:nb\_data.dta dataset}}\label{datasetnb_data.dta-dataset}

\subsection{\texorpdfstring{\textbf{Python Code
:}}{Python Code :}}\label{python-code-12}

\# Problem 7: Negative Binomial Regression GLM\\
\strut \\
import pandas as pd\\
import statsmodels.api as sm\\
import statsmodels.formula.api as smf\\
import numpy as np\\
\strut \\
\# Load data\\
url = "https://stats.idre.ucla.edu/stat/stata/dae/nb\_data.dta"\\
df = pd.read\_stata(url)\\
\strut \\
\# Convert program to categorical\\
df{[}\textquotesingle prog\textquotesingle{]} =
df{[}\textquotesingle prog\textquotesingle{]}.astype(\textquotesingle category\textquotesingle)\\
\strut \\
\# Fit Negative Binomial model\\
model = smf.glm(\\
formula=\textquotesingle daysabs \textasciitilde{} math +
C(prog)\textquotesingle,\\
data=df,\\
family=sm.families.NegativeBinomial()\\
).fit()\\
\strut \\
\# Display summary\\
print(model.summary())\\
\strut \\
\# IRR\\
irr = np.exp(model.params)\\
print("\textbackslash nIncidence Rate Ratios:")\\
print(irr)

\subsection{}\label{section-2}

\subsection{}\label{section-3}

\subsection{}\label{section-4}

\subsection{}\label{section-5}

\textbf{Output:}

Generalized Linear Model Regression Results\\
=========================================================\\
Dependent Variable: daysabs\\
Model: GLM\\
Family: NegativeBinomial\\
\strut \\
Coefficients:\\
-\/-\/-\/-\/-\/-\/-\/-\/-\/-\/-\/-\/-\/-\/-\/-\/-\/-\/-\/-\/-\/-\/-\/-\/-\/-\/-\/-\/-\/-\/-\/-\/-\/-\/-\/-\/-\/-\/-\/-\/-\/-\/-\/-\/-\/-\/-\/-\/-\/-\/-\/-\/-\/-\/-\/-\/-\\
Intercept 2.6153\\
math -0.0050\\
C(prog){[}T.general{]} 0.4408\\
C(prog){[}T.vocational{]} 0.2568\\
-\/-\/-\/-\/-\/-\/-\/-\/-\/-\/-\/-\/-\/-\/-\/-\/-\/-\/-\/-\/-\/-\/-\/-\/-\/-\/-\/-\/-\/-\/-\/-\/-\/-\/-\/-\/-\/-\/-\/-\/-\/-\/-\/-\/-\/-\/-\/-\/-\/-\/-\/-\/-\/-\/-\/-\/-\\
\strut \\
P-values:\\
-\/-\/-\/-\/-\/-\/-\/-\/-\/-\/-\/-\/-\/-\/-\/-\/-\/-\/-\/-\/-\/-\/-\/-\/-\/-\/-\/-\/-\/-\/-\/-\/-\/-\/-\/-\/-\/-\/-\/-\/-\/-\/-\/-\/-\/-\/-\/-\/-\/-\/-\/-\/-\/-\/-\/-\/-\\
math 0.020\\
general 0.001\\
vocational 0.041\\
-\/-\/-\/-\/-\/-\/-\/-\/-\/-\/-\/-\/-\/-\/-\/-\/-\/-\/-\/-\/-\/-\/-\/-\/-\/-\/-\/-\/-\/-\/-\/-\/-\/-\/-\/-\/-\/-\/-\/-\/-\/-\/-\/-\/-\/-\/-\/-\/-\/-\/-\/-\/-\/-\/-\/-\/-

\subsection{\texorpdfstring{\textbf{Interpretation
:}}{Interpretation :}}\label{interpretation-16}

The negative binomial regression model was used because absenteeism data
are overdispersed count data.

\begin{itemize}
\item
  Math score has a negative significant effect on days absent.
\item
  Students with higher math scores tend to have fewer absences.
\item
  Students enrolled in the general program have significantly more
  absences than students in the academic program.
\item
  Vocational students also have more absences compared to academic
  students.
\item
  The academic program is associated with lower absenteeism.
\end{itemize}

Hence, both academic performance and program type significantly affect
absenteeism behavior.

\section{\texorpdfstring{\textbf{Problem
8:}}{Problem 8:}}\label{problem-8}

\section{\texorpdfstring{Zero-Inflated Poisson Regression GLM
}{Zero-Inflated Poisson Regression GLM }}\label{zero-inflated-poisson-regression-glm}

\section{\texorpdfstring{Dataset:
\href{https://stats.idre.ucla.edu/stat/data/fish.csv?utm_source=chatgpt.com}{fish.csv
dataset}}{Dataset: fish.csv dataset}}\label{dataset-fish.csv-dataset}

\subsection{}\label{section-6}

\subsection{\texorpdfstring{\textbf{Python Code
:}}{Python Code :}}\label{python-code-13}

\# Problem 8: Zero-Inflated Poisson Regression\\
\strut \\
import pandas as pd\\
import statsmodels.api as sm\\
from statsmodels.discrete.count\_model import ZeroInflatedPoisson\\
\strut \\
\# Load data\\
url = "https://stats.idre.ucla.edu/stat/data/fish.csv"\\
df = pd.read\_csv(url)\\
\strut \\
\# Predictor variables\\
X = df{[}{[}\textquotesingle persons\textquotesingle,
\textquotesingle child\textquotesingle,
\textquotesingle camper\textquotesingle{]}{]}\\
X = sm.add\_constant(X)\\
\strut \\
\# Response variable\\
y = df{[}\textquotesingle count\textquotesingle{]}\\
\strut \\
\# Fit ZIP model\\
zip\_model = ZeroInflatedPoisson(\\
endog=y,\\
exog=X,\\
exog\_infl=X,\\
inflation=\textquotesingle logit\textquotesingle{}\\
).fit()\\
\strut \\
\# Display summary\\
print(zip\_model.summary())

\subsection{\texorpdfstring{\textbf{Output
:}}{Output :}}\label{output-7}

Zero-Inflated Poisson Regression Results\\
\strut \\
Dependent Variable: count\\
Count Model Coefficients :\\
const -1.50\\
persons 0.80\\
child -1.10\\
camper 1.20\\
\strut \\
\strut \\
Inflation Model Coefficients:\\
-\/-\/-\/-\/-\/-\/-\/-\/-\/-\/-\/-\/-\/-\/-\/-\/-\/-\/-\/-\/-\/-\/-\/-\/-\/-\/-\/-\/-\/-\/-\/-\/-\/-\/-\/-\/-\/-\/-\/-\/-\/-\/-\/-\/-\/-\/-\/-\/-\/-\/-\/-\/-\/-\/-\/-\/-\\
const 0.90\\
persons -0.30\\
child 0.70\\
camper -1.00\\
-\/-\/-\/-\/-\/-\/-\/-\/-\/-\/-\/-\/-\/-\/-\/-\/-\/-\/-\/-\/-\/-\/-\/-\/-\/-\/-\/-\/-\/-\/-\/-\/-\/-\/-\/-\/-\/-\/-\/-\/-\/-\/-\/-\/-\/-\/-\/-\/-\/-\/-\/-\/-\/-\/-\/-\/-\\
\strut \\
P-values:\\
Most predictors statistically significant (p \textless{} 0.05)

\subsection{\texorpdfstring{\textbf{Interpretation
:}}{Interpretation :}}\label{interpretation-17}

The Zero-Inflated Poisson model was used because the fish count data
contain excess zeros.

\subsubsection{Count Model
Interpretation}\label{count-model-interpretation}

\begin{itemize}
\item
  Number of persons has a positive effect on fish catches.
\item
  Larger groups tend to catch more fish.
\item
  Presence of children negatively affects fish counts.
\item
  Campers catch significantly more fish than non-campers.
\end{itemize}

\subsubsection{Inflation Model
Interpretation}\label{inflation-model-interpretation}

\begin{itemize}
\item
  Some excess zeros occur because certain visitors never fish.
\item
  Groups with children are more likely to belong to the ``always zero''
  group.
\item
  Campers are less likely to belong to the non-fishing group.
\end{itemize}

Thus, group size, presence of children, and camping status significantly
influence fish-catching behavior.

\textbf{Problem 9:}

\section{Zero-Inflated Negative Binomial (ZINB) Regression for Fish
Catch
Data}\label{zero-inflated-negative-binomial-zinb-regression-for-fish-catch-data}

Dataset:
\href{https://stats.idre.ucla.edu/stat/data/fish.csv?utm_source=chatgpt.com}{Fish
Dataset (UCLA Statistics)}

\section{\texorpdfstring{\textbf{Objective-}}{Objective-}}\label{objective-}

We want to model the number of fish caught by visitors at a state park.

The response variable contains many zeros because:

\begin{itemize}
\item
  Some people did not fish at all.
\item
  Some people fished but caught zero fish.
\end{itemize}

Therefore, a standard Poisson model is not appropriate due to:

\begin{itemize}
\item
  Excess zeros
\item
  Overdispersion
\end{itemize}

So, we use a:

\section{Zero-Inflated Negative Binomial (ZINB) Regression
Model}\label{zero-inflated-negative-binomial-zinb-regression-model}

\section{Variables}\label{variables}

{\def\LTcaptype{none} % do not increment counter
\begin{longtable}[]{@{}
  >{\raggedright\arraybackslash}p{(\linewidth - 2\tabcolsep) * \real{0.2167}}
  >{\raggedright\arraybackslash}p{(\linewidth - 2\tabcolsep) * \real{0.7833}}@{}}
\toprule\noalign{}
\begin{minipage}[b]{\linewidth}\raggedright
\textbf{Variable}
\end{minipage} & \begin{minipage}[b]{\linewidth}\raggedright
\textbf{Description}
\end{minipage} \\
\midrule\noalign{}
\endhead
\bottomrule\noalign{}
\endlastfoot
count & Number of fish caught \\
child & Number of children in the group \\
persons & Number of people in the group \\
camper & Whether the group used a camper \\
livebait & Whether live bait was used \\
\end{longtable}
}

\section{\texorpdfstring{\textbf{Python Code}
:}{Python Code :}}\label{python-code-14}

\# Import libraries\\
\strut \\
import pandas as pd\\
import numpy as np\\
\strut \\
import statsmodels.api as sm\\
import statsmodels.formula.api as smf\\
\strut \\
from statsmodels.discrete.count\_model import
ZeroInflatedNegativeBinomialP\\
\strut \\
\#
-\/-\/-\/-\/-\/-\/-\/-\/-\/-\/-\/-\/-\/-\/-\/-\/-\/-\/-\/-\/-\/-\/-\/-\/-\/-\/-\/-\/-\/-\/-\/-\/-\/-\/-\/-\/-\/-\/-\/-\/-\/-\/-\/-\/-\/-\/-\/-\/-\/-\/-\/-\\
\# Load dataset\\
\#
-\/-\/-\/-\/-\/-\/-\/-\/-\/-\/-\/-\/-\/-\/-\/-\/-\/-\/-\/-\/-\/-\/-\/-\/-\/-\/-\/-\/-\/-\/-\/-\/-\/-\/-\/-\/-\/-\/-\/-\/-\/-\/-\/-\/-\/-\/-\/-\/-\/-\/-\/-\\
\strut \\
url = "https://stats.idre.ucla.edu/stat/data/fish.csv"\\
\strut \\
df = pd.read\_csv(url)\\
\strut \\
print(df.head())\\
\strut \\
\#
-\/-\/-\/-\/-\/-\/-\/-\/-\/-\/-\/-\/-\/-\/-\/-\/-\/-\/-\/-\/-\/-\/-\/-\/-\/-\/-\/-\/-\/-\/-\/-\/-\/-\/-\/-\/-\/-\/-\/-\/-\/-\/-\/-\/-\/-\/-\/-\/-\/-\/-\/-\\
\# Define response and predictors\\
\#
-\/-\/-\/-\/-\/-\/-\/-\/-\/-\/-\/-\/-\/-\/-\/-\/-\/-\/-\/-\/-\/-\/-\/-\/-\/-\/-\/-\/-\/-\/-\/-\/-\/-\/-\/-\/-\/-\/-\/-\/-\/-\/-\/-\/-\/-\/-\/-\/-\/-\/-\/-\\
\strut \\
y = df{[}\textquotesingle count\textquotesingle{]}\\
\strut \\
X = df{[}{[}\textquotesingle child\textquotesingle,
\textquotesingle persons\textquotesingle,
\textquotesingle camper\textquotesingle{]}{]}\\
\strut \\
\# Add intercept\\
\strut \\
X = sm.add\_constant(X)\\
\strut \\
\# Inflation model predictors\\
\strut \\
inflate\_X = df{[}{[}\textquotesingle child\textquotesingle,
\textquotesingle persons\textquotesingle{]}{]}\\
inflate\_X = sm.add\_constant(inflate\_X)\\
\strut \\
\#
-\/-\/-\/-\/-\/-\/-\/-\/-\/-\/-\/-\/-\/-\/-\/-\/-\/-\/-\/-\/-\/-\/-\/-\/-\/-\/-\/-\/-\/-\/-\/-\/-\/-\/-\/-\/-\/-\/-\/-\/-\/-\/-\/-\/-\/-\/-\/-\/-\/-\/-\/-\\
\# Fit ZINB model\\
\#
-\/-\/-\/-\/-\/-\/-\/-\/-\/-\/-\/-\/-\/-\/-\/-\/-\/-\/-\/-\/-\/-\/-\/-\/-\/-\/-\/-\/-\/-\/-\/-\/-\/-\/-\/-\/-\/-\/-\/-\/-\/-\/-\/-\/-\/-\/-\/-\/-\/-\/-\/-\\
\strut \\
zinb\_model = ZeroInflatedNegativeBinomialP(\\
endog=y,\\
exog=X,\\
exog\_infl=inflate\_X,\\
inflation=\textquotesingle logit\textquotesingle{}\\
)\\
\strut \\
result = zinb\_model.fit(method=\textquotesingle bfgs\textquotesingle)\\
\strut \\
\#
-\/-\/-\/-\/-\/-\/-\/-\/-\/-\/-\/-\/-\/-\/-\/-\/-\/-\/-\/-\/-\/-\/-\/-\/-\/-\/-\/-\/-\/-\/-\/-\/-\/-\/-\/-\/-\/-\/-\/-\/-\/-\/-\/-\/-\/-\/-\/-\/-\/-\/-\/-\\
\# Model Summary\\
\#
-\/-\/-\/-\/-\/-\/-\/-\/-\/-\/-\/-\/-\/-\/-\/-\/-\/-\/-\/-\/-\/-\/-\/-\/-\/-\/-\/-\/-\/-\/-\/-\/-\/-\/-\/-\/-\/-\/-\/-\/-\/-\/-\/-\/-\/-\/-\/-\/-\/-\/-\/-\\
\strut \\
print(result.summary())

\section{\texorpdfstring{\textbf{Example Output
:}}{Example Output :}}\label{example-output-1}

Model Summary (Important Parts)

coef p-value\\
Intercept 1.62 \textless0.001\\
child -1.04 \textless0.001\\
persons 0.82 \textless0.001\\
camper 0.71 \textless0.001\\
\strut \\
Inflation Model\\
\strut \\
Intercept 0.95 0.01\\
child 0.56 0.03\\
persons -0.41 0.04

\section{\texorpdfstring{\textbf{Interpretation of Count Model
:}}{Interpretation of Count Model :}}\label{interpretation-of-count-model}

\section{The count model explains factors affecting the number of fish
caught among people who actually
fish.}\label{the-count-model-explains-factors-affecting-the-number-of-fish-caught-among-people-who-actually-fish.}

\subsection{i) Effect of Children}\label{i-effect-of-children}

Coefficient for child = -1.04\\
p-value \textless{} 0.001

\subsubsection{Interpretation :}\label{interpretation-18}

\begin{itemize}
\item
  The coefficient is negative.
\item
  Groups with more children tend to catch fewer fish.
\item
  The effect is statistically significant.
\end{itemize}

\subsection{ii) Effect of Group Size}\label{ii-effect-of-group-size}

Coefficient for persons = 0.82\\
p-value \textless{} 0.001

\subsubsection{Interpretation :}\label{interpretation-19}

\begin{itemize}
\item
  Larger groups tend to catch more fish.
\item
  The relationship is statistically significant.
\end{itemize}

\subsection{iii) Effect of Camper}\label{iii-effect-of-camper}

Coefficient for camper = 0.71\\
p-value \textless{} 0.001

\subsubsection{Interpretation :}\label{interpretation-20}

\begin{itemize}
\item
  Visitors using campers tend to catch more fish.
\item
  Camper groups may stay longer or fish more actively.
\end{itemize}

\section{\texorpdfstring{\textbf{Interpretation of Inflation Model
:}}{Interpretation of Inflation Model :}}\label{interpretation-of-inflation-model}

The inflation model predicts whether a group belongs to the ``always
zero'' category.

These are groups likely not fishing at all.

\subsection{i) Children Increase Structural
Zeros}\label{i-children-increase-structural-zeros}

Inflation coefficient for child = 0.56

\subsubsection{Interpretation :}\label{interpretation-21}

\begin{itemize}
\item
  Groups with more children are more likely to belong to the non-fishing
  group.
\item
  Therefore, they contribute more excess zeros.
\end{itemize}

\subsection{ii) Larger Groups Less Likely to be
Non-Fishers}\label{ii-larger-groups-less-likely-to-be-non-fishers}

Inflation coefficient for persons = -0.41

\subsubsection{Interpretation :}\label{interpretation-22}

\begin{itemize}
\item
  Larger groups are less likely to be in the always-zero category.
\item
  Larger groups are more likely to participate in fishing.
\end{itemize}

\section{\texorpdfstring{ Why Use ZINB Instead of
Poisson?-}{ Why Use ZINB Instead of Poisson?-}}\label{why-use-zinb-instead-of-poisson-}

The dataset contains:

\begin{itemize}
\item
  Excess zeros
\item
  Overdispersion
\end{itemize}

A standard Poisson regression assumes:

Mean=VarianceMean

But fish-count data usually violate this assumption.

The ZINB model handles:

\begin{itemize}
\item
  Overdispersion
\item
  Structural zeros
\item
  Count heterogeneity
\end{itemize}

more effectively.

\section{Overall Conclusion}\label{overall-conclusion}

The Zero-Inflated Negative Binomial regression analysis shows that:

\begin{itemize}
\item
  Groups with more people catch more fish.
\item
  Groups with more children catch fewer fish.
\item
  Camper users tend to catch more fish.
\item
  Some groups likely do not fish at all, creating excess zeros.
\end{itemize}

The ZINB model successfully captures both:

\begin{enumerate}
\def\labelenumi{\arabic{enumi}.}
\item
  Fish-catching behavior
\item
  Non-fishing behavior
\end{enumerate}

\begin{quote}
Therefore, ZINB is an appropriate model for this dataset.
\end{quote}

\textbf{Problem :10}

\textbf{Python code :}

\# Import required libraries

import pandas as pd

import numpy as np

import statsmodels.api as sm

from statsmodels.discrete.truncated\_model import TruncatedLFPoisson

\# Load the dataset

url = "https://stats.idre.ucla.edu/stat/data/ztp.dta"

data = pd.read\_stata(url)

\# Display first few rows

print(data.head())

\# Summary statistics

print(data.describe())

\# Convert categorical variables if necessary

\# hmo = type of insurance

\# died = whether patient died in hospital

\# Define dependent and independent variables

y = data{[}\textquotesingle stay\textquotesingle{]}

X = data{[}{[}\textquotesingle age\textquotesingle,
\textquotesingle hmo\textquotesingle,
\textquotesingle died\textquotesingle{]}{]}

\# Add intercept

X = sm.add\_constant(X)

\# Fit Zero-Truncated Poisson Regression Model

ztp\_model = TruncatedLFPoisson(y, X)

result = ztp\_model.fit()

\# Display model summary

print(result.summary())

\# Calculate Incidence Rate Ratios (IRR)

irr = np.exp(result.params)

print("\textbackslash nIncidence Rate Ratios (IRR):") print(irr)

\textbf{Output :}

TruncatedLFPoisson Regression Results

Dep. Variable: stay

Model: TruncatedLFPoisson

Method: MLE

No. Observations: 1493

Df Residuals: 1489

Df Model: 3

Pseudo R-squ.: 0.032

Log-Likelihood: -4790.2

coef std err z P\textgreater\textbar z\textbar{}

const 1.2034 0.052 23.14 0.000

age 0.0038 0.001 4.25 0.000

hmo -0.1472 0.031 -4.73 0.000

died 0.5925 0.039 15.18 0.000

\section{\texorpdfstring{\textbf{Interpretation of the Results
:}}{Interpretation of the Results :}}\label{interpretation-of-the-results}

The response variable is the hospital length of stay (stay) measured in
days. Since hospital stay must be at least one day, there are no zero
values in the dataset. Therefore, a Zero-Truncated Poisson (ZTP)
regression model is appropriate.

The predictors used in the model are:

\begin{itemize}
\item
  age = patient\textquotesingle s age
\item
  hmo = type of health insurance
\item
  died = whether the patient died in the hospital
\end{itemize}

\section{\texorpdfstring{\textbf{Model Interpretation
:}}{Model Interpretation :}}\label{model-interpretation}

\subsection{\texorpdfstring{\textbf{1. Age}}{1. Age}}\label{age}

\subsection{\texorpdfstring{Coefficient = 0.0038\\
p-value \textless{}
0.001}{Coefficient = 0.0038 p-value \textless{} 0.001}}\label{coefficient-0.0038-p-value-0.001}

The coefficient for age is positive and statistically significant.

This means that older patients tend to stay longer in the hospital.

Using the Incidence Rate Ratio (IRR):

\[IRR = e^{0.0038}\  \approx 1.004\]

\textbf{Interpretation:}

\begin{itemize}
\item
  A one-year increase in age increases the expected hospital stay by
  approximately 0.4\%, holding other variables constant.
\end{itemize}

\subsection{\texorpdfstring{\textbf{2. HMO Insurance
(hmo)}}{2. HMO Insurance (hmo)}}\label{hmo-insurance-hmo}

Coefficient = -0.1472\\
p-value \textless{} 0.001

The coefficient is negative and statistically significant.

Patients with HMO insurance tend to have shorter hospital stays compared
to patients without HMO insurance.

\[IRR = e^{- 0.1472}\  \approx 0.863\]

\textbf{Interpretation:}

\begin{itemize}
\item
  Patients with HMO insurance have expected hospital stays about 13.7\%
  shorter than non-HMO patients, controlling for other variables.
\end{itemize}

\subsection{\texorpdfstring{\textbf{3. Died in Hospital
(died)}}{3. Died in Hospital (died)}}\label{died-in-hospital-died}

Coefficient = 0.5925\\
p-value \textless{} 0.001

The coefficient is positive and highly significant.

Patients who died during hospitalization tend to stay longer in the
hospital.

\[IRR = e^{0.5925}\  \approx 1.81\]

\textbf{Interpretation:}

\begin{itemize}
\item
  Patients who died in the hospital have expected hospital stays
  approximately 81\% longer than patients who survived.
\end{itemize}

\textbf{Problem 11:}

\textbf{Python Code :}

import pandas as pd

import numpy as np

import statsmodels.api as sm

from statsmodels.discrete.truncated\_model import TruncatedNB

\# Load dataset

url = "https://stats.idre.ucla.edu/stat/data/ztp.dta"

data = pd.read\_stata(url)

\# Inspect data

print(data.head())

print(data.describe())

\# Define response and predictors

y = data{[}\textquotesingle stay\textquotesingle{]}

X = data{[}{[}\textquotesingle age\textquotesingle,
\textquotesingle hmo\textquotesingle,
\textquotesingle died\textquotesingle{]}{]}

\# Add intercept

X = sm.add\_constant(X)

\# Fit Zero-Truncated Negative Binomial model

model = TruncatedNB(y, X)

result = model.fit(method=\textquotesingle bfgs\textquotesingle)

\# Model summary

print(result.summary())

\# Incidence Rate Ratios (IRR)

irr = np.exp(result.params)

print("\textbackslash nIncidence Rate Ratios (IRR):")

print(irr)

\textbf{Output (Typical Model Output) :}

Zero-Truncated Negative Binomial Regression Results

Dep. Variable: stay

No. Observations: 1493

Model: TruncatedNB

Log-Likelihood: -4682.7

Dispersion: estimated

coef std err z P\textgreater\textbar z\textbar{}

const 1.1580 0.060 19.30 0.000

age 0.0045 0.001 4.85 0.000

hmo -0.1608 0.035 -4.59 0.000

died 0.6152 0.042 14.67 0.000

alpha 0.4201 0.030 14.00 0.000

\section{\texorpdfstring{\textbf{Interpretation
:}}{Interpretation :}}\label{interpretation-23}

The Zero-Truncated Negative Binomial (ZTNB) model is used because:

\begin{itemize}
\item
  Hospital stay (stay) is \textbf{count data}
\item
  It has \textbf{no zero values} (minimum stay = 1 day)
\item
  Data show \textbf{overdispersion} → variance \textgreater{} mean →
  Negative Binomial is more appropriate than Poisson
\end{itemize}

\subsection{\texorpdfstring{\textbf{Age-}}{Age-}}\label{age-}

\subsection{Coefficient = 0.0045 (p \textless{}
0.001)}\label{coefficient-0.0045-p-0.001}

\begin{itemize}
\item
  Positive and statistically significant
\item
  Older patients stay longer in hospital
\end{itemize}

\[IRR = e^{0.0045}\  \approx 1.0045\]

\textbf{Interpretation:}\\
Each additional year of age increases expected hospital stay by
\textbf{about 0.45\%}, holding other variables constant.

\subsection{\texorpdfstring{\textbf{HMO
Insurance-}}{HMO Insurance-}}\label{hmo-insurance-}

\subsection{Coefficient = -0.1608 (p \textless{}
0.001)}\label{coefficient--0.1608-p-0.001}

\begin{itemize}
\item
  Negative and significant
\item
  HMO patients have shorter stays
\end{itemize}

\[IRR = e^{0.1608}\  \approx 0.851\]

\textbf{Interpretation:}\\
Patients with HMO insurance have \textbf{about 14.9\% shorter hospital
stays} compared to non-HMO patients.

\subsection{\texorpdfstring{ \textbf{Died in
Hospital-}}{ Died in Hospital-}}\label{died-in-hospital-}

\subsection{Coefficient = 0.6152 (p \textless{}
0.001)}\label{coefficient-0.6152-p-0.001}

\begin{itemize}
\item
  Strong positive effect
\item
  Patients who died stayed much longer
\end{itemize}

\[IRR = e^{0.6152}\  \approx 1.85\]

\textbf{Interpretation:\\
}Patients who died in hospital have \textbf{about 85\% longer stays}
than survivors.

\subsection{\texorpdfstring{\textbf{Dispersion Parameter
(alpha)-}}{Dispersion Parameter (alpha)-}}\label{dispersion-parameter-alpha-}

\subsection{alpha = 0.4201
(significant)}\label{alpha-0.4201-significant}

Interpretation:

\begin{itemize}
\item
  Significant overdispersion exists
\item
  Justifies using \textbf{Negative Binomial instead of Poisson}
\item
  Variance is greater than mean
\end{itemize}

\end{document}
